\section{Findings}\label{sec:findings}

\subsection{Weights decide which jobs go first}\label{sec:weights}

\sumpara On the hand-checkable case, the whole gain comes from moving a low-weight job behind two high-weight jobs. \sumparaend
\fin \Cref{fig:gantt} draws both schedules of the 6-job instance. \fine
\fin Earliest due date places job J005, due at 67 with weight 1, ahead of jobs J003 and J004, which carry weight 4. \fine
\fin Iterated greedy moves J005 to the end, so J003 finishes at 97 and J004 at 119, against 108 and 139 under the rule, while J005 finishes at 136, against 82. \fine
\fin The weighted tardiness of J003 and J004 falls by 124 weighted minutes and that of J005 rises by 54, and with a smaller rise for J006 the total falls from 281 to 233. \fine
\fin The same reordering raises waiting on this instance from 13 to 37~minutes, because the two jobs now at the end, J006 and J005, wait 18 and 19~minutes at stages S2 to S4 behind the long operations of the weight-4 jobs. \fine

\begin{figure}[ht]
\centering
{\footnotesize\color{muted}Earliest due date: J001, J002, J005, J003, J006, J004 (total weighted tardiness 281)}\\[2pt]
\gantt{\eddschedule}\\[6pt]
{\footnotesize\color{muted}Iterated greedy: J001, J002, J003, J004, J006, J005 (total weighted tardiness 233, the optimum)}\\[2pt]
\gantt{\igschedule}\\[4pt]
{\footnotesize\color{muted}Shading gives the job's weight:
\textcolor{w1}{\rule{0.7em}{0.7em}}~1\quad
\textcolor{w2}{\rule{0.7em}{0.7em}}~2\quad
\textcolor{w4}{\rule{0.7em}{0.7em}}~4}
\caption{Both schedules of the 6-job instance on stages S1 to S4, with time in minutes on the horizontal axis. Each bar is one operation, labeled with its job number and shaded by the job's weight. Iterated greedy moves the weight-1 job J005 behind the weight-4 jobs J003 and J004.}
\label{fig:gantt}
\end{figure}

\sumpara On the development instances, the same weighting turns most late jobs into on-time jobs. \sumparaend
\fin On the 45-job instance, the share of jobs finishing on time rose from 0.244 under earliest due date to 0.911 under iterated greedy, and on the 15-job and 120-job instances from 0.333 to 0.867 and from 0.233 to 0.642. \fine

\subsection{The search removes most of the tardiness the weighted start leaves, and 120 jobs need more of it}\label{sec:search}

\sumpara The weighted start already beats the rule, and the search removes most of what remains. \sumparaend
\fin NEH\_EWDD alone lowered weighted tardiness relative to the rule by 95.1, 84.8, and 58.9~percent on the 15-, 45-, and 120-job instances (\cref{tab:comparison}). \fine
\fin The iterations of \cref{alg:ig} then lowered NEH\_EWDD's value by a further 70.0, 83.9, and 53.6~percent, or 49, 302, and 5,069 weighted minutes. \fine
\fin On the 120-job instance the search had not converged after 5000~iterations, since five other seeds found sequences 1.8 to 5.1~percent below the default seed's 4,383. \fine
\fin Better sequences are therefore within reach of the same search on the largest instance, at a computation time that already reaches 37.2~seconds. \fine

\sumpara Breaking ties by waiting matters once tardiness reaches its floor. \sumparaend
\fin On the 15-job instance, the solver's optimal order and iterated greedy's order both have tardiness 21, but the solver's has 323~minutes of waiting and iterated greedy's 252 (\cref{tab:comparison}). \fine
\fin Among sequences with optimal tardiness, the tie-break thus removes 71~minutes of work in process that a tardiness-only objective leaves in place. \fine

\subsection{Variability makes the searched plan's on-time jobs late, without thinner margins}\label{sec:mechanism}

\sumpara Under variability the searched plan loses on-time jobs, while the rule's plan mainly lengthens lateness it already has. \sumparaend
\fin At $\nu = 0.2$, the iterated greedy plan had 4 late jobs as planned and 11.5 on average across the 30~scenarios, while the rule's plan had 34 and 36.1. \fine
\fin The rule's mean weighted tardiness rose by 1,170.1 weighted minutes over its plan while its count of late jobs rose by only 2.1, so most of the increase came from jobs that were already late finishing later still. \fine
\fin The iterated greedy plan's mean rose by 362.6, a factor of 7.25 over its planned 58, against a factor of 1.49 for the rule, because its losses come from jobs that were on time in the plan. \fine

\sumpara The searched plan's on-time margins are as wide as the rule's; it is exposed because it has more on-time jobs. \sumparaend
\fin The on-time jobs of the iterated greedy plan had a median slack of 31~minutes before their due dates, against 28 for the rule's plan. \fine
\fin In each plan, 3 on-time jobs had slack below one standard deviation of their own total processing time at $\nu = 0.2$, which is $0.2\sqrt{\sum_k p_{jk}^2}$ for job $j$. \fine
\fin The iterated greedy plan is exposed because it has 41 on-time jobs that a delay can make late, against 11 for the rule. \fine
